% 第21版：法稿物理78--79页及80页NB，英编稿51页中部--52页NB末。
% 与 initial 模块接续；此文件不含13.4标题。
\source{法文手稿物理第78--79页及第80页顶部附记；英文编订稿第51页中部至第52页附记末}

现在只剩下证明引理13.3。按通常的做法，只须证明：经过$x$、余维为$n-1$的一般线性子簇，记作$L$，具有性质a)、b)、c)。性质a)，以及b)、c)中的维数结论（对于任意给定的射影浸入），都立即由2.3得到；这里2.3已在第12节重新审视并作了修正。像第8节中已经用过的一次论证那样，将它应用于参数化经过$x$的直线的射影空间，以及$Z$经以$x$为中心的锥投影在该空间中的像。\cmt{V21-BER001}{法稿第78页“余维”中的\emph{co}、b)、c)的维数内容以及“对于任意给定的射影浸入”均为可辨的作者手添语，正文恢复最后这一英文遗漏的限定。法稿两处都写小写$x$，英稿其中一处误排大写$X$；旧D17-M062的改正结论保留，但归责改为英文转写错误。这里经过$x$的条件在法稿证明中明确存在；本模块仍须与前面引理陈述及有理点条件一并阅读。}

［此外，把用这种方法得到的某些结果明确写出，或许也有益处：这些结果涉及由满足经过某个固定线性子簇这一条件的线性子簇所截出的线性截面。可以在正文中写出，也可以另设一节。］\cmt{V21-BER002}{方括号及编写建议属于原作者，不是译者新加的一项定理。旧D17-M063的身份说明继续保留。}

对于b)、c)中的光滑性，由于a)，可以分别用$X\setminus Z$、$Y\setminus Z$代替$X$、$Y$；它们都是光滑的。于是归结为证明以下断言。

设$f\colon X\to P$是一个无分歧态射，$X$在$k$上光滑，并取$x\in P(k)$，使得$x$不属于$X$中任何维数小于$m$的分量的像。\cmt{V21-BER003}{法稿写小写$x\in P(k)$，没有英稿附加的“不可约分量？”这一问句；后者作为英译编者的疑问保存于此。光滑概形的局部诺特正规性使不可约分量局部开闭；本次独立补证明列域上有限型条件，以保证只须处理有限多个分量，不把原段的简写扩大为任意无有限性条件的预概形。}
令$\eta$为$\operatorname{Grass}^{m}(P)$中由经过$x$的余维$m$线性子簇$L$组成的子格拉斯曼簇的一般点。则相应的一般线性截面$\underline X^{(m)}_\eta$在$\kappa(\eta)$上光滑；至少当$k$的特征为零时如此，而在另一种情形中，须将$f$换成$\varphi_d\circ f$，其中$d$为任意不小于2的整数。\cmt{V21-BER004}{法稿已经明确写出带下划线的$\underline X$及下标$\eta$；英稿在本句漏去一般点下标。法稿和英稿都把光滑性的基域写为$k$，正文订为一般纤维的定义域$\kappa(\eta)$。这是数学订正，不能误说法稿本来就写出了$\kappa(\eta)$：例如$\Spec k(t)\to\Spec k$并非局部有限表示的光滑态射。法稿同页两次均为$\varphi_n f$，英稿先排$Q_nf$、后排$\Phi_nf$；此项法英记号差异现已查明。正文将Veronese次数改记为$d$，以区别于下面表示$\dim X$的$n$；“任意”在法稿中清楚可读。报告V21-BER-E03给出$d\geq2$、有限型范围内的独立一阶截面补证，不以此声称原8.8缺文已恢复。}

这是应补入第9节的一项内容，而它本身又来自应补入第8节的下述内容。沿用8.8的记号，并假设$X$不可约；对于我们正在处理的问题，这样做是允许的。\cmt{V21-BER005}{法稿清楚写\emph{Pb}，即“问题”的缩写；英稿误作$p^0$并加问号。旧D17-M067所保留的这个未辨记号现可撤销。\emph{remords}在此仍表示对先前章节的补写建议。原8.8完整起首及8.7(c)的有关缺续并未因此恢复，其依赖边界仍保留。}
若
\[
\operatorname{codim}T\geq2,
\]
或者$\underline Y^{\mathrm{sing}}\to T$一般平展（这一条件在$k$的特征为零时自动满足；或者在将$f$换成$\varphi_d\circ f$、$d\geq2$之后满足，参见第9节），则对于经过$x$的一般超平面$H_\eta$，截面$\underline X^{(1)}_\eta$光滑，维数为$n-1$；须除去$f(X)=\{x\}$的情形，此时$n=0$。\cmt{V21-BER006}{法稿此处打字为带下划线的$Y^{\mathrm{sing}}$，下一页则为带下划线的$X^{\mathrm{sing}}$；英稿在这里误转为$h^{\mathrm{sing}}$。正文按各页保留法稿记号，不悄悄将二者统一，更不能把它们解释为当前光滑$X$自身的奇异点集。法稿写的是固定$x$、超平面一般，以及$f(X)=\{x\}$；英稿的“a generic x”和$f(x)$均属转写或语序问题。关于“一般平展”的括注是原作者对先前结果的引用，其一般性并未因文字恢复而获全面独立证明。}

暂且承认这一结果，它显然解决了上述补写内容中$m=1$的特殊情形。注意，允许作基域扩张后，$\underline X^{(m)}_\eta$可以如下得到：取一个经过$x$、余维为$m-1$的$L'$，以及一个经过$x$、余维为1的$H$，使$L'$和$H$对于这些条件是相互独立的一般选择（用通常的严谨说法，就是置于参数化二元组$(L',H)$的概形的一般点处），再取$X$由$L'\cap H$给出的线性截面。可以先取由$H$给出的截面；由前面的结果，它是光滑的。然后再由$L'$截取；由归纳假设，所得截面光滑。因此，对$m$作归纳，便得到一般情形。\cmt{V21-BER007}{法稿第79页的两步及其理由完整可读：先由$H$截取，用刚承认的$m=1$结果；再由$L'$截取，用归纳假设。英稿漏掉$L'\cap H$及关键句，旧D17-M070、M071曾根据残缺英文整理为先$L'$后$H$，本版恢复法稿顺序，并撤销“法稿语法不完整”这一归责。这里的基域扩张和二元组参数空间均为原文已有内容。}

这种类型的论证，此前已经用于把2.6等结果推广到任意余维$m$的线性截面。值得在一般情形下一次性把它清楚写出，以便以后引用，而不必每次都进入一份正式表述中那些略显繁重的细节。\cmt{V21-BER008}{法稿末句\emph{des détails un peu lourds d'une présentation en forme}完整可辨。英稿的法语标记、空方括号和连写\emph{enforme}不再当作原法稿缺文；旧D17-M072相应撤销。此处仍是作者的编写建议。}

还须证明上面宣布的8.8的推论，即$m=1$的情形。如果$\operatorname{codim}T\geq2$，另一方面，参数化满足$x\in H_\xi$的超平面的子空间$Q$是对偶射影空间$P^\vee$中的一个超平面，余维为1；因此它的一般点$\eta$不可能属于$T$，所需结论随即得到。\cmt{V21-BER009}{法稿\emph{on gagne}是完整的结句；英稿另附“已经完成？”的翻译疑问，不属于原作者的数学疑问。维数论证使用$T$为闭集：若$Q$的一般点属于$T$，则$Q\subset T$，与$T$余维至少2冲突。}

在$\operatorname{codim}T=1$的情形，由于$T$不可约，只有当$Q=T$时，才可能有$\eta\in T$。用几何语言来说（可以假设$k$代数闭），这意味着：对每个$z\in X$，$X$在$z$处的切空间，确切地说，是该切空间经$f'_z$所得的像，经过$x$。\cmt{V21-BER010}{法稿明确写\emph{que si $Q=T$}，即“只有当$Q=T$时”；英稿的\emph{under}造成含义不明，不能归咎于原法稿。法稿关于代数闭域的括注也有明确范围，英稿缺失了括号结构。切空间的说法在这里是射影几何中的切空间像，不是说抽象切向量空间含有一个射影点。}
我们来证明，在$\underline X^{\mathrm{sing}}\to T$一般平展、也就是处于所引用的8.8条件之下时，这种情形不能发生，除非$f(X)=\{x\}$，此时$X$为零维。\cmt{V21-BER011}{法稿两处清楚写\emph{lorsque}，即“当……时”；英文却两次写\emph{unless}，把所需条件的逻辑方向译反。旧D17-M076正文采用的方向正确，但本版据法稿确认这是恢复原意，而非替原作者修补逻辑。法稿此页确实改写为带下划线的$X^{\mathrm{sing}}$，与前页$Y^{\mathrm{sing}}$的区别仍记录；在有限型无分歧背景下，像为一点才使定义域为零维，不能将此推论外推到任意常值态射。}

事实上，由8.7(c)（它本质上表达了$X$与其“对偶”$T$之间关系的对称性），对于几乎所有$z\in X(k)$，$f(z)$与$T$在某一点处的切空间正交。由于$T=Q$，这就意味着$f(z)$与$Q$正交，故$f(z)=x$，从而$f(X)=\{x\}$。这就完成了上述补写内容的证明，从而也完成了13.2的证明。\cmt{V21-BER012}{法稿的\emph{donc}、手添\emph{orthogonal}以及两次\emph{d'où}均可辨，推论顺序完整；英稿相应语序破碎。这里“几乎所有”是一般点意义，而非测度意义。原作者所说“完成”仍依赖其引用的8.7(c)、8.8；本版分别记录这一作者结语、尚存缺续，以及V21-BER-E03在明确范围内另外给出的证明，不混为同一件事。}

\noindent\textbf{附记。}
如果把$x$换成一个维数大于0的线性子簇$C$，并要求$H$经过$C$，上述论证就不能照样进行。事实上，没有什么能排除$T$包含直线$C^0$的情形（例如，取$\dim C=r-2$；这是仍使$H$能够实际变化的最大维数）；可以取$P$中的一个非奇异二次超曲面作为例子，特征任意。\cmt{V21-BER013}{法稿第80页\emph{rien ne s'oppose ... que T contienne la droite $C^0$}完整清楚。英文的“没有理由认为”“without supposing (?)”等混乱不再作为法稿残句保留，旧D17-M079的主要疑点现已消除。这里$C^0$按射影对偶语境表示所有包含$C$的超平面的参数子空间；当$\dim C=r-2$时它是一条直线。法稿首个子簇字母覆打，后文明确用$C$，正文据同段统一。冒号后有两行密集删划，已识别为作者删稿；其完整旧句不能可靠复原，源页图像保留，不以未删文字倒填。报告给出光滑二次曲面的一族全坏平面截面的直接核算。}

\noindent\textbf{译校订正。}即使将$f$换成$\varphi_d\circ f$、$d\geq2$，若对固定的正维线性子簇$C$不另加限制，也仍可能出现这样的现象；因此不能据此断言经过$C$的一般截面光滑。\cmt{V21-BER014}{原作者最后一句完整意思是：“不过，对于$\varphi_n f$、$n\geq2$，这样的现象或许不再发生；我们可以把这个问题作为第9节的一条注提出。”它是一项疑问或猜测，不是作者声称已证明的定理。其不加限制的肯定答案有反例：从$f=\id_{\mathbf P^1}$出发，作三次Veronese嵌入$[u:v]\mapsto[u^3:u^2v:uv^2:v^3]$；取其在$[1:0:0:0]$处的切线$C=V(x_2,x_3)$。包含$C$的平面束$a x_2+b x_3=0$拉回为$v^2(au+bv)=0$，每个成员都不光滑，一般成员具有一个二重点及另一个单点。此$C$与像相交；反例不涉及另加$C\cap(\varphi_d\circ f)(X)=\varnothing$等条件的版本。原问句及旧D17-M080的未定身份继续在本注与报告中保存。}
% 精确止点：法80 NB末“au n°9”；下一为 Remarques13.4(a)。
